\section{Results}

\subsection{Overview and evidence tiers}
\label{sec:results-overview}

Results follow the three evidence tiers of Table~\ref{tab:objective-evidence-map}.
\textbf{Part~A} (\S\ref{sec:results-assembly-preprocessing}) is primary evidence on implementation auditability and hydrography preprocessing (Objectives~1--2); \textbf{Part~B} (\S\ref{sec:results-cal-val-sensitivity}) is secondary evidence on structural scaling and assembly cost (Objective~3) and downstream workflow readiness (Objective~4); \textbf{Part~C} (\S\ref{sec:results-runtime-benchmark}) is supporting evidence on simulation runtime (Objective~5).

\subsection{Assembly and preprocessing evidence}
\label{sec:results-assembly-preprocessing}

\subsubsection{Objective 1: Implementation and workflow architecture}
\label{sec:results-obj1}

The automated workflow produces a standard SWAT\textsuperscript{+} project tree---vector hydrography, gridded and station climate inputs, SQLite project tables, and optional QA sidecars---through the layered pipeline of Figure~\ref{fig:workflow}, with asynchronous workers (Figure~\ref{fig:architecture-layers}) separating long-running assembly from interactive use so outputs remain in a predictable, auditable per-basin layout (representative file classes in Table~\ref{tab:package-manifest-abridged}; full path patterns in Supplementary Table~S1). Structural footprints for the S/M/L size classes are in Table~\ref{tab:benchmark-sm-l-products} (\S\ref{sec:results-obj3}).

\subsubsection{Objective 2: Hydrography preprocessing outcomes}
\label{sec:results-obj2}

National NHDPlus HR vectors are transformed into SWAT\textsuperscript{+}-compatible channel, subbasin, and lake layers through the staged pipeline of Figure~\ref{fig:nhd-workflow}; Table~\ref{tab:nhd-rules-abridged} lists the core rules and Table~\ref{tab:nhd-preprocessing-qa} compares the original USGS NHDPlus HR (HU4 geodatabase, domain-clipped) with the exported shapefiles for the S/M/L domains, documenting where divergence-code-2, isolated, and orphan/no-VAA catchments are consolidated (e.g.\ 56$\rightarrow$53 watershed polygons for S; 9{,}739 original catchments to 162 subbasins and 8{,}181 channels for L). These are structural-rule diagnostics, not demonstrated hydrologic benefit.

\subsubsection{Objective 2 (continued): Drainage-area fidelity audit}
\label{sec:results-drainage-audit}

Structural counts confirm the rules were applied but not that the \emph{executable} model preserves cumulative drainage structure. To close that gap, \texttt{chandeg.con} channel area was compared with original NHDPlus HR \texttt{TotDASqKm} at gage-assigned channels (\S\ref{sec:methods-drainage-audit}) across all eight Model IDs (Table~\ref{tab:drainage-area-audit}). Of 98 audited gages, 97 matched a SWAT\textsuperscript{+} channel and 82 (85\%) fell within the loose $0.5$--$2.0\times$ band. The band is deliberately wide: a gage can snap to a tributary, canal, or lake-outlet reach whose cumulative area differs sharply from the intended mainstem, so a factor-of-two band isolates true structural mismatch from discrete assignment error. Where assignment is least ambiguous---the gage-dense Peace River domain---the distribution tightens well inside the band: 66 of 76 gages agree, median ratio $1.11$, $10$th--$90$th percentile $0.78$--$1.28$ (central 80\% within $\approx\pm25\%$). Per-model medians span $0.57$--$1.16$, with the lowest values driven by a few assignment outliers on the sparsely gaged scaling domains rather than basin-wide conversion error.

\begin{table}[htbp]
  \centering
  \caption{Drainage-area fidelity audit: SWAT\textsuperscript{+} \texttt{chandeg.con} area versus original NHDPlus HR \texttt{TotDASqKm} at gage-assigned channels, by catalog Model ID (Objective~2). ``Outliers'' are gage--channel assignment outliers reported separately from the agreement band.}
  \label{tab:drainage-area-audit}
  \footnotesize
  \setlength{\tabcolsep}{4pt}
  \renewcommand{\arraystretch}{1.05}
  \resizebox{\textwidth}{!}{%
    % Generated by publication/analysis/scripts/emit_tab_drainage_area_audit_tex.py
% Source: /data/SWATGenXApp/codes/web_application/frontend/src/data/drainageAreaAuditCatalog.json (lastUpdated 2026-06-02)
\begin{tabular}{@{}lllrrrrr@{}}
\toprule
Tier & Catalog ID & State & Stations & Matched & Median SWAT/NHD & Within 0.5--2.0$\times$ & Outliers \\
\midrule
S & \texttt{03080102} & FL & 3 & 3 & 0.85 & 3 & 0 \\
M & \texttt{09471300} & AZ & 3 & 3 & 1.16 & 3 & 0 \\
L & \texttt{03100101} & FL & 76 & 76 & 1.11 & 66 & 6 \\
X20 & \texttt{03152000} & WV & 6 & 6 & 0.84 & 4 & 2 \\
X40 & \texttt{07174000} & KS & 3 & 2 & 1.00 & 2 & 0 \\
X60 & \texttt{15060105} & AZ & 4 & 4 & 0.57 & 2 & 2 \\
Cal & \texttt{02297600} & FL & 1 & 1 & 1.08 & 1 & 0 \\
Cal & \texttt{05536265} & IL & 2 & 2 & 0.88 & 1 & 1 \\
\midrule
\textbf{All} & --- & --- & \textbf{98} & \textbf{97} & --- & \textbf{82} & \textbf{11} \\
\bottomrule
\end{tabular}
%
  }
\end{table}

\paragraph{Station--channel assignment and agreement (Peace River).}
The whole-HUC-8 Peace domain (\texttt{03100101}; 76 gages) supports per-class analysis. The NHD-first/SWAT-second method (Table~\ref{tab:station-assignment}) mapped 61 gages by exact NHD$\rightarrow$GIS crosswalk, 10 by lake-outlet replacement, and 5 by controlled downstream replacement, yielding 36 clean tributaries, 2 clean mainstem reaches, 16 mainstem reaches with a known NHD offset, and 22 lake/canal cases flagged for review; a reconciliation funnel left 16 fully calibration-ready gages (54 passed the assignment class, 20 with usable 2000--2024 daily streamflow). All 76 gages matched a SWAT\textsuperscript{+} channel (median ratio 1.11, 66 inside the loose band; Figure~\ref{fig:drainage-area-audit-scatter}). The NHD reference also tracks the independent USGS NWIS catalog (Table~\ref{tab:drainage-credibility}): of 57 gages with a site area, 43 agree within 25\% (median NHD/NWIS 0.85), and the two alignment questions are near-independent ($r=-0.005$), so catalog variance does not drive the model--hydrography offset (per-class detail in Supplementary Section~\ref{sec:supp-delineation}).

\begin{table}[htbp]
  \centering
  \caption{NHD-first/SWAT-second USGS station-to-channel assignment for the Peace River HUC-8 domain (Objective~2).}
  \label{tab:station-assignment}
  \small
  \setlength{\tabcolsep}{6pt}
  \renewcommand{\arraystretch}{1.05}
  % Generated by publication/analysis/scripts/emit_tab_station_assignment_tex.py
% Source: /data/SWATGenXApp/codes/web_application/frontend/src/data/stationAssignmentPageCatalog.json
\begin{tabular}{@{}lr@{}}
\toprule
Peace River station assignment (76 USGS gages) & Gages \\
\midrule
\multicolumn{2}{@{}l}{\emph{Assignment class}} \\
\quad Clean tributary & 36 \\
\quad Mainstem, clean & 2 \\
\quad Mainstem, known NHD offset & 16 \\
\quad Lake-outlet review & 13 \\
\quad Canal/artificial review & 9 \\
\midrule
\multicolumn{2}{@{}l}{\emph{Mapping method}} \\
\quad Exact NHD$\rightarrow$GIS crosswalk & 61 \\
\quad Lake-outlet replacement & 10 \\
\quad Controlled downstream replacement & 5 \\
\midrule
\multicolumn{2}{@{}l}{\emph{Calibration reconciliation funnel}} \\
\quad Auto-calibration assignment class & 54 \\
\quad \dots\ with usable 2000--2024 daily flow ($\geq$80\%) & 20 \\
\quad Fully calibration-ready & 16 \\
\bottomrule
\end{tabular}
%
\end{table}

\begin{table}[htbp]
  \centering
  \caption{Peace River two-way drainage-area alignment: NHD reference versus USGS NWIS site area, and joint NHD--USGS / SWAT\textsuperscript{+}--NHD agreement (Objective~2). The two comparisons are near-independent ($r=-0.005$).}
  \label{tab:drainage-credibility}
  \small
  \setlength{\tabcolsep}{6pt}
  \renewcommand{\arraystretch}{1.05}
  % Generated by publication/analysis/scripts/emit_tab_drainage_credibility_tex.py
% Source: /data/SWATGenXApp/codes/web_application/frontend/src/data/stationAssignmentPageCatalog.json
\begin{tabular}{@{}lr@{}}
\toprule
Drainage-area alignment, Peace River & Count / ratio \\
\midrule
\multicolumn{2}{@{}l}{\emph{NHD reference vs USGS NWIS site area} (n=57 with site area; 19 lack it)} \\
\quad NHD within 10\% of NWIS & 15 \\
\quad NHD within 25\% of NWIS & 43 \\
\quad NHD $>$25\% from NWIS & 14 \\
\quad Median NHD/NWIS ratio & 0.85 \\
\midrule
\multicolumn{2}{@{}l}{\emph{Joint NHD--USGS and SWAT--NHD agreement} (n=57)} \\
\quad Both agree ($|$NHD$-$USGS$|\leq$25\%, SWAT/NHD 0.5--2.0$\times$) & 40 \\
\quad NHD--USGS mismatch only & 13 \\
\quad SWAT/NHD mismatch only & 3 \\
\quad Both mismatch & 1 \\
\bottomrule
\end{tabular}
%
\end{table}

\begin{table}[htbp]
  \centering
  \caption{Peace River mainstem drainage-area pipeline trace on a diagnostic panel of representative mainstem reaches: median SWAT\textsuperscript{+}/NHD ratio by stage, localizing the systematic offset (Objective~2). Stage labels defined in Section~\ref{sec:methods-drainage-audit}.}
  \label{tab:drainage-area-pipeline-trace}
  \footnotesize
  \setlength{\tabcolsep}{4pt}
  \renewcommand{\arraystretch}{1.1}
  \resizebox{\textwidth}{!}{%
    % Generated by publication/analysis/scripts/emit_tab_drainage_area_pipeline_trace_tex.py
% Source: /data/SWATGenXApp/codes/publication/analysis/qa/peace-drainage-area-pipeline-trace.csv (mainstem diagnostic-panel gages; n=3)
\begin{tabular}{@{}clcp{5.0cm}@{}}
\toprule
Stage & Quantity & Median ratio vs NHD \texttt{TotDASqKm} & Interpretation \\
\midrule
A & Original NHDPlus HR HU4 \texttt{TotDASqKm} (VAA) & 1.00 & Reference cumulative drainage at the gaged reach \\
B & Original upstream $\Sigma$ \texttt{AreaSqKm} (HU4 network) & 1.00 & Independent cumulative cross-check; agrees with A \\
G & Cleaned flowline reach \texttt{TotDASqKm} & 1.00 & Pre-QSWAT\textsuperscript{+} attribute; $<$0.1\% vs A (VAA not pre-inflated) \\
C & Cleaned flowline-network upstream $\Sigma$ & 0.99 & After topology cleaning; no positive bias \\
D & Post-processed watershed-polygon upstream $\Sigma$ & 1.03 & SWATGenX subbasin polygons; $+2$--$3\%$ only \\
E & QSWAT\textsuperscript{+}/TauDEM channel \texttt{AreaC} at assignment & 1.17 & First large exceedance ($+15$--$17\%$) \\
F & \texttt{chandeg.con} executable area & 1.17 & Headline audit column; matches E \\
H & Live \texttt{rivs1.AreaC} / project SQLite & 1.17 & Matches E/F ($<$0.1\% vs \texttt{chandeg}); offset present before export \\
\bottomrule
\end{tabular}
%
  }
\end{table}

\begin{figure}[htbp]
  \centering
  \includegraphics[width=0.72\textwidth]{final/fig-drainage-area-audit-peace-scatter.png}
  \caption{Peace River HUC-8 drainage-area agreement: SWAT\textsuperscript{+} \texttt{chandeg.con} area versus original NHDPlus HR \texttt{TotDASqKm} at 76 gage-assigned channels (log--log), colored by SWAT\textsuperscript{+}/NHD agreement band; the $1{:}1$ line and the systematic mainstem $+10$--$17\%$ offset are annotated.}
  \label{fig:drainage-area-audit-scatter}
\end{figure}

\paragraph{Where the offset enters.}
The 16 mainstem-offset gages carry a systematic $+10$--$17\%$ excess in SWAT\textsuperscript{+} cumulative area. Tracing that area through the pipeline (Table~\ref{tab:drainage-area-pipeline-trace}) localizes it: the original VAA, upstream catchment-area sums, the pre-QSWAT\textsuperscript{+} cleaned-flowline attribute, and the post-processed watershed polygons all agree with the reference within $\pm3\%$ (median ratios $0.99$--$1.03$), while the excess first appears at the QSWAT\textsuperscript{+}/TauDEM channel-area stage (\texttt{AreaC}, ratio $\approx1.17$) and persists into \texttt{chandeg.con}. The live project channel area (\texttt{rivs1.AreaC} and SQLite \texttt{gis\_channels}) carries the full $+17\%$ and \texttt{chandeg.con} matches it to within $0.1\%$. The NHDPlus HR preprocessing therefore preserves cumulative drainage structure; the residual is a delineation-definition difference between TauDEM channel contributing area and NHD hydrosequence cumulative drainage, introduced downstream of the preprocessing this paper documents. These are structural-fidelity diagnostics, not hydrologic validation.

\subsubsection{Objective 3: Structural scalability and model-assembly cost}
\label{sec:results-obj3}

The small, medium, and large benchmark domains (Table~\ref{tab:model-roster}; \texttt{03080102} Oklawaha FL, \texttt{09471300} Upper San Pedro AZ, \texttt{03100101} Peace River FL) span 473 to 94{,}303 HRUs (Figure~\ref{fig:example-basin-maps}). Table~\ref{tab:benchmark-sm-l-products} combines discretization counts, on-disk size, and measured model-assembly wall time on the production host (end-to-end worker job for S/M; maximum logged QSWAT\textsuperscript{+} assembly phase for the whole-HUC-8 L domain); simulation runtime is reported separately in Part~C.

\begin{figure}[htbp]
  \centering
  \includegraphics[width=\textwidth]{final/fig-example-basin-maps-combined-3panel.png}
  \caption{Small, medium, and large benchmark domains (size classes S, M, L). Panels (a)--(c) are independently zoomed and are not at a common geographic scale. Structure and model-assembly cost: Table~\ref{tab:benchmark-sm-l-products}.}
  \label{fig:example-basin-maps}
\end{figure}

\subsection{Downstream workflow readiness}
\label{sec:results-cal-val-sensitivity}

\subsubsection{Objective 4: Workflow execution on controlled gage basins}
\label{sec:results-obj4}

Two controlled gage basins (\texttt{02297600} Florida, \texttt{05536265} Illinois), separate from the S/M/L set, test whether generated projects enter a standard downstream workflow (\S\ref{sec:methods-calibration}). On both, the platform completed initialization pools, PSO calibration, independent-period forward runs, and Morris batch execution with ensemble diagnostics on non-overlapping scored windows (Florida holdout 2019--2024; Illinois holdout 2012--2015; Table~\ref{tab:controlled-gage-structure}). These results do not isolate individual preprocessing effects or establish national validity.

\begin{table}[htbp]
  \centering
  \caption{Structural context for Objective~4 controlled gage basins.}
  \label{tab:controlled-gage-structure}
  \scriptsize
  \setlength{\tabcolsep}{3pt}
  \renewcommand{\arraystretch}{1.05}
  \resizebox{\textwidth}{!}{%
    % Generated by publication/analysis/scripts/emit_tab_controlled_gage_structure_tex.py
\begin{tabular}{@{}llrrrrrrrrllr@{}}
\toprule
Basin ID & State & Ch. & Orig.\ cat. & Orig.\ reach & Div-2 & Isol.\ reach & Final ch. & Final sub. & Cal.\ period & Ver.\ period & HRUs \\
\midrule
\texttt{02297600} & FL & 2 & 41 & 41 & 0 & 4 & 37 & 4 & 2013-2018 & 2019-2024 & --- \\
\texttt{05536265} & IL & 25 & 224 & 224 & 6 & 13 & 206 & 8 & 2020-2024 & 2012-2015 & --- \\
\bottomrule
\end{tabular}
%
  }
\end{table}

\paragraph{Morris sensitivity and ensemble uncertainty.}
Morris analysis (1000 evaluations; 27 trajectories, four levels) on each basin's calibration window yields a 999-member 2.5--97.5\% predictive band (Figure~\ref{fig:sensitivity-ensemble}; Table~\ref{tab:sensitivity-ensemble-metrics}); P-factor and R-factor characterize that band, not overall goodness-of-fit. Florida's ensemble is narrower (daily P-factor~0.71, R-factor~0.73) than Illinois (0.87, 3.13). Ranked $\mu^*$ (Tables~\ref{tab:morris-sensitivity}, \ref{tab:morris-sensitivity-05536265}) differ by climate: Florida is driven by \texttt{surq\_lag}, \texttt{perco}, \texttt{alpha\_bf}, and \texttt{cn3\_swf} with negligible snow terms, whereas Illinois is driven by snow (\texttt{melt\_min}, \texttt{melt\_max}) and channel routing (\texttt{k}, \texttt{mann}). That driver set \emph{explains} the weaker Illinois fit and wider band---snow forcing and routing are degenerate against streamflow alone---which is the intended use of these diagnostics: to surface where and why a generated model is hard to calibrate, not to certify skill (spider charts and monthly ensembles in Supplementary Section~\ref{sec:supp-calval-sensitivity}).

\begin{table}[htbp]
  \centering
  \caption{Morris ensemble uncertainty diagnostics (999 completed members; P-factor and R-factor defined in Methods).}
  \label{tab:sensitivity-ensemble-metrics}
  \footnotesize
  \setlength{\tabcolsep}{4pt}
  \renewcommand{\arraystretch}{1.05}
  % Generated by emit_tab_sensitivity_ensemble_metrics_tex.py
% Source: /data/SWATGenXApp/codes/publication/tables/tab-sensitivity-ensemble-metrics.csv
\begin{tabular}{@{}llrrrr@{}}
\toprule
Basin & Step & $N$ & P-factor & R-factor & Best NSE \\
\midrule
\texttt{02297600} & Daily & 999 & 0.705 & 0.733 & 0.817 \\
\texttt{02297600} & Monthly & 999 & 0.722 & 0.894 & 0.882 \\
\texttt{05536265} & Daily & 999 & 0.874 & 3.125 & 0.290 \\
\texttt{05536265} & Monthly & 999 & 0.900 & 5.123 & 0.470 \\
\bottomrule
\end{tabular}
%
\end{table}

\begin{figure}[htbp]
  \centering
  \includegraphics[width=\textwidth,keepaspectratio]{final/fig-sensitivity-ensemble-controlled-basins-daily.png}
  \caption{Morris ensemble daily streamflow for controlled gage basins on each basin's calibration window: (a)~\texttt{02297600} (Florida) and (b)~\texttt{05536265} (Illinois). Shaded band: 2.5--97.5\% of 999 Morris members; dashed gray: ensemble median; navy: observed; dark red: best Morris member. P-factor and R-factor annotated on each panel.}
  \label{fig:sensitivity-ensemble}
\end{figure}

\begin{table}[htbp]
  \centering
  \caption{Top eight Morris $\mu^*$ for basin \texttt{02297600} (full 36-parameter ranking in archived \texttt{tab-sensitivity-morris.csv}).}
  \label{tab:morris-sensitivity}
  \footnotesize
  \setlength{\tabcolsep}{4pt}
  \renewcommand{\arraystretch}{1.05}
  % Generated by emit_tab_sensitivity_morris_tex.py
% Source: /data/SWATGenXApp/codes/publication/tables/tab-sensitivity-morris.csv
% site_no: 02297600
\begin{tabular}{@{}rlrrr@{}}
\toprule
Rank & Parameter & $\mu$ & $\mu^*$ & $\sigma$ \\
\midrule
1 & \texttt{surq\_lag} & -0.6435 & 0.6435 & 0.4685 \\
2 & \texttt{perco} & 0.2323 & 0.2880 & 0.4113 \\
3 & \texttt{cn3\_swf} & 0.0598 & 0.1413 & 0.1699 \\
4 & \texttt{alpha\_bf} & -0.0921 & 0.0921 & 0.0987 \\
5 & \texttt{dep\_wt} & 0.0609 & 0.0882 & 0.1565 \\
6 & \texttt{spec\_yld} & 0.0523 & 0.0545 & 0.0945 \\
7 & \texttt{flo\_min} & -0.0142 & 0.0519 & 0.0774 \\
8 & \texttt{dp\_es} & 0.0500 & 0.0500 & 0.0282 \\
\bottomrule
\end{tabular}
%
\end{table}

\begin{table}[htbp]
  \centering
  \caption{Top eight Morris $\mu^*$ for basin \texttt{05536265} (full 36-parameter ranking in archived \texttt{tab-sensitivity-morris.csv}).}
  \label{tab:morris-sensitivity-05536265}
  \footnotesize
  \setlength{\tabcolsep}{4pt}
  \renewcommand{\arraystretch}{1.05}
  \input{tables/tab-sensitivity-morris-05536265}%
\end{table}

\paragraph{Split-sample calibration diagnostics.}
Full daily/monthly NSE, RMSE, and PBIAS by stage are in Supplementary Table~\ref{tab:hydrologic-metrics}; holdout skill is summarized here and is not interpreted as platform validation.
On \textbf{Florida} (\texttt{02297600}, channel~2; calibration 2013--2018 after a three-year warm-up, daily NWIS from 2009-10-01), the initialization pool best captured seasonal timing (daily NSE~0.80), calibration reached daily/monthly NSE~0.84/0.90, and independent verification retained 0.73/0.80 with modest wet bias (PBIAS~$\approx+7\%$).
On \textbf{Illinois} (\texttt{05536265}, channel~25; calibration 2020--2024, pre-calibration holdout 2012--2015), skill was lower throughout (calibration daily/monthly NSE~0.37/0.60; verification 0.22/0.42), reported as a workflow diagnostic on a second hydroclimate---automated assembly does not ensure calibration-ready hydrology, and Morris still completed (999 members) with the wider band noted above.
\label{sec:results-cal-val-05536265}

\subsection{Operational simulation cost}
\label{sec:results-runtime-benchmark}

\subsubsection{Objective 5: SWAT\textsuperscript{+} simulation runtime benchmark}
\label{sec:results-obj5}

Objective~5 reports measured simulation cost on one documented host/build (AMD EPYC 7282, 10~vCPU; Intel ifx \texttt{-O3 -ipo}; calendar-2021 365-day runs under the filtered NetCDF \texttt{channel\_sd} profile). Filtered runs scale from about 5~s (473 HRUs) to about 100~s (11{,}284 HRUs) and about 23~min (94{,}303 HRUs) per simulated year (Table~\ref{tab:runtime-benchmark}; HRU scaling in Figure~\ref{fig:runtime-benchmark-hru-scaling}). Print-scope, export-format, compiler, and the six-point X20/X40/X60 ladder appear in Supplementary Section~\ref{sec:supp-runtime-benchmark} and on the public benchmark page \citep{SWATGenXcompanion}.

\begin{table}[htbp]
  \centering
  \caption{One-year simulation wall time under calibration filtered NetCDF profile for the small, medium, and large benchmark domains.}
  \label{tab:runtime-benchmark}
  \small
  % Generated by emit_tab_runtime_benchmark_tex.py
% Source: /data/SWATGenXApp/codes/publication/tables/tab-runtime-benchmark.csv
\begin{tabular*}{\textwidth}{@{\extracolsep{\fill}}llrrrr@{}}
\toprule
Size class & Model & HRUs & Wall (s) & s/day & Peak RSS (KB) \\
\midrule
S & \texttt{03080102} & 473 & 5.42 & 0.01 & 55924 \\
M & \texttt{09471300} & 11284 & 99.79 & 0.26 & 511468 \\
L & \texttt{03100101} & 94303 & 1378.13 & 3.37 & 3956648 \\
\bottomrule
\end{tabular*}
%
\end{table}

\begin{figure}[htbp]
  \centering
  \includegraphics[width=0.78\textwidth]{final/fig-runtime-benchmark-hru-scaling.png}
  \caption{Wall time versus HRU count for the six published benchmark domains (Supplementary Table~\ref{tab:runtime-benchmark-hru-scaling}).}
  \label{fig:runtime-benchmark-hru-scaling}
\end{figure}
